% GENERATED by CasDesignBasis/tools/design_report.py from design_basis.toml v2.3.2. Do not edit.
\begin{tabularx}{\textwidth}{@{}l r l l L@{}}
\toprule
Symbol & Value & Unit & Source & Reason \\
\midrule
$\dot m_\text{max}$ & \num{104.2} & kg/h & [DB] tab. 13, [CT] & Maximum O$_2$ gas flow per oxygenator (CT10, sea water, 15 \% gas/liquid ratio). An oxygenator adds O$_2$ only while its own inlet pump runs. \\
$\eta$ & \num{0.8} & -- & [CTE] & Transfer efficiency; the test gave 0.81 initial efficiency, rounded down.\assumed{} \\
$T_f$ & \num{2} & s & design choice & Closed-loop time constant of FIC41-43, used in the control design model. \\
$\tau_v$ & \num{0.22} & s & [B2865], [S418] & Open-loop lag of valve and flow meter, used in the simulator: valve $<25$ ms, flow meter $t_{90}=0.5$ s ($\tau_v = t_{90}/\ln 10$). \\
$u_\text{cut}$ & \num{2} & \% & [B8605] & Driver cut-off; below this the valve is closed (Linde default). Above it the gas flow is linear in the command up to $\dot m_\text{max}$. \\
$u_\%$ & 4-20 mA & -- & user, [B8605] & 4 mA (0 \%) = fully closed, 20 mA (100 \%) = fully open. The PLC writes 0--32760 for 0--100 \% to the 750-563 output; the Bürkert 8605 driver takes the 4--20 mA (Linde default) and sets the coil current. Below the cut-off the valve stays closed, so a lost signal is expected to close the valve (no O$_2$ on that line). \\
$\tau_s$ & \num{19.54} & s & [RDO] & O$_2$ sensor as first order fitted to $T_{90}<45$ s ($\tau_s=T_{90}/\ln 10$). [RDO] gives $T_{63}<5$ s, so this is conservative. \\
$c_\text{sat}$ & 9.09 & g/m$^3$ & [O2S] & Saturation at $T_w$, $S_w$ (Weiss). \\
$c_\text{in}$ & 7.27 & g/m$^3$ &  & O$_2$ in the intake water. \\
$c_\text{ref}$ & 8.18 & g/m$^3$ &  & O$_2$ setpoint. \\
$\dot m_\text{max}$ & 28.94 & g/s &  & Maximum O$_2$ gas flow per line. \\
\bottomrule
\end{tabularx}
